The slow depolarizing state \(u\) (V2.1) was the designated candidate
for a persistent, intrinsic memory that does not depend on synaptic
structure. Phase I's evidence against that role is unusually complete
because the audits measured what \(u\) actually does during and after
the stimulus, and because one claim was retracted when its runs turned
out not to exist.

\subsection{What \(u\) is, measured}

\(u\) tracks the organism's own spiking through a first-order filter:
per-spike increment \(\beta\), decay \(\tau_s\). Three measured facts
\cite{infoaudit}:

\begin{enumerate}
\item \textbf{It saturates during the stimulus.} Means rise
monotonically with cumulative exposure and plateau at a flux balance
(W1: \(\bar u\to 1.32\) by presentation \(\sim\)30; W2:
\(\bar u\to 0.86\), still rising at drive end). A saturating integrator
cannot encode cumulative episode history beyond its equilibrium.
\item \textbf{Persistence without structure.} After stimulus removal
the population does not relax: the active core self-sustains as
deterministic pacemaking --- 500 spikes/s on the refractory ceiling
(2 ms refractory \(\Rightarrow\) 500 Hz), Fano factors \(\approx 0\),
5--6 active neurons of 52, Gini concentration \(0.80 \to 0.93\) over
silence. The persistent activity is a clock, not a code
\cite{infoaudit}.
\item \textbf{The only tested claim of a discriminative \(u\) is
retracted.} The V2.1 Stage-C validation table claimed temporal-capacity
unchanged at the rule-selected operating point; the referenced runs do
not exist in the repository. The information audit discovered the
phantom reference, retracted the verdict as ungrounded (``unsupported,
not refuted --- unmeasured''), and added an integrity regression that
fails any record citing nonexistent runs \cite{infoaudit}.
\end{enumerate}

\subsection{The historical wedge: \(u\) fails during the stimulus, not only after it}

The standard memory objection to an integrator is that it decays after
stimulus offset. The ANIMA record shows something stronger: in the
regimes where episode identity was tested, the persistent quantities
fail to distinguish episodes \emph{while the stimulus is still
present}. The alternating-store cosine after adjacent \(A\) and \(C\)
presentations converges to 0.99 (\S\ref{sec:bottleneck}) --- the state
after an \(A\) presentation is already near-identical to the state after
the adjacent \(C\), at a measurement taken 2 s apart with exactly one
intervening presentation \cite{synmem}. The V2.2 latch design
(discussed next) was in part motivated by this timing: it attempted to
make a per-neuron bistable bit the carrier, under u-driven SET/RESET
thresholds --- and failed as a carrier (latches decode the curriculum at
chance in all 120 stage-2 runs; \(\eta_{\mathrm{rel}}=1.0\)
self-cancels) \cite{overview}. The E18--E21 temporal-capacity series
bracketed the same wedge from the output side: no antecedent-dependent
divergence over 800 ms; output responses stimulus-locked with zero
endogenous output activity; temporal capacity below 50 ms
\cite{overview}.

\subsection{The drive-gated slow state and the writable-memory review}

The X-series explored the one amendment to \(u\) that its own audit
motivated: gate the per-spike write by the neuron's afferent drive
(\texttt{slow\_state\_beta\_drive}; drive trace
\(x_i = \min(I_{\mathrm{aff}}/v_{\mathrm{th}}, 1)\), fitted with
the frozen membrane scale after the dimensional correction of
\S\ref{sec:repro}; spec and logs
\cite{xspecdrive,xplorlog}). The records measure:
\begin{itemize}
\item X1: a transition band at \(\beta \approx 0.0047\) --- an
intermediate regime with a decaying tail whose effective time constant
(\(\approx\) 8.2 s) differs from \(\tau_u\);
\item X3: a 45-run boundary map in which top-\(u\) at drive end is a
near-perfect predictor of the post-drive regime (\(\rho \approx
0.97\)); drive composition shifts regime, but seed dominates;
\item the architectural review \cite{xplor2}: sensory organization
exists at the afferent layer (42--48 of 52 neurons dominant for a given
pure-drive curriculum; interleaving keeps separable cohorts);
during-stimulus responses separate (cosine 0.30--0.89) while
off-window responses do not (0.999+); the persistent state is a
``seed-lottery'' twice over.
\end{itemize}
The review's primary diagnosis --- \emph{insufficient writable /
competitive memory}; the missing capability is coexisting
independently writeable traces under budget competition --- is the
exact bridge to the synaptic-store chapters: it reframes the negative
as a memory-capability question and names budget competition (M2) as
the suspected mechanism, which the SDE probes and the CLLA era then
isolated \cite{xplor2,overview}.

\begin{quote}
\textbf{Demonstrated:} scalar intrinsic persistence (\(u\)) is
self-sustaining but information-poor (saturated mean, pacemaker core,
Fano \(\approx 0\)); the V2.2 latch carrier failed 120-run tests; the
temporal carry of the substrate is below 50 ms.\\[1pt]
\textbf{Supported:} episode discriminability fails inside the stimulus
window, hence the bottleneck is not post-offset decay but
representation-integration collapse.\\[1pt]
\textbf{Not established / retracted:} the Stage-C capacity claim is
ungrounded and retracted; no A/C contrast in \(u\) was ever executed.
\end{quote}